\section*{Resumen}
\addcontentsline{toc}{section}{Resumen}

Este proyecto compara tres colas de prioridad dentro del mismo algoritmo de Dijkstra para resolver el problema de caminos mínimos desde un origen único en redes viales urbanas: un min-heap binario indexado, un heap 4-ario indexado y un Fibonacci heap. El caso aplicado utiliza una consulta sobre una red vial real de Cusco obtenida de OpenStreetMap (891 vértices y 2139 aristas dirigidas), mientras que la comparación controlada se realiza sobre instancias sintéticas reproducibles con perfiles disperso ($|E|=4|V|$) y denso ($|E|=0.08\,|V|(|V|-1)$). El protocolo mide tiempo con \texttt{perf\_counter\_ns}, desactiva el recolector de basura durante cada muestra, ejecuta 15 repeticiones sobre tres instancias por escala y reporta mediana e intervalo intercuartílico; la corrección se valida contra NetworkX y un certificado de optimalidad de Bellman en $O(|V|+|E|)$.

Los resultados no permiten declarar una cola ganadora universal. El heap 4-ario obtiene la menor mediana en 7 de 10 escalas sintéticas, con una razón mediana de 0.96 frente al heap binario en el perfil disperso y 0.95 en el denso, pero los intervalos intercuartílicos se solapan en las 10 escalas y el heap binario gana en 3 escalas y en la consulta OSM (15.65 ms frente a 19.87 ms del 4-ario y 26.09 ms del Fibonacci). El Fibonacci heap nunca obtiene la menor mediana y resulta entre 1.05 y 1.59 veces más lento que el binario. La extensión $A^*$ reduce las expansiones de 887 a 823 ($-7.2\%$) sin alterar el costo óptimo. Un factor meteorológico uniforme de 1.0691 incrementa el tiempo estimado en $6.91\%$, pero no modifica la ruta porque escala todos los pesos por igual. La conclusión central es que la elección de cola depende de la escala, la densidad y el entorno de ejecución, y que la ventaja asintótica del Fibonacci heap no se traduce en ventaja práctica en esta implementación de Python.

\vspace{0.4cm}
\noindent\textbf{Palabras clave:} algoritmo de Dijkstra, colas de prioridad, heap binario, heap 4-ario, Fibonacci heap, rutas mínimas, OpenStreetMap, Cusco, $A^*$, análisis experimental.

\vspace{0.8cm}
\section*{Abstract}
\addcontentsline{toc}{section}{Abstract}

This project compares three priority queues within the same Dijkstra algorithm for the single-source shortest path problem on urban road networks: an indexed binary min-heap, an indexed 4-ary heap, and a Fibonacci heap. The applied case uses a real road query in Cusco retrieved from OpenStreetMap (891 vertices and 2139 directed edges), while the controlled comparison uses reproducible synthetic instances with sparse ($|E|=4|V|$) and dense ($|E|=0.08\,|V|(|V|-1)$) profiles. The protocol measures wall-clock time with \texttt{perf\_counter\_ns}, disables garbage collection during each sample, runs 15 repetitions on three instances per scale, and reports the median and interquartile range; correctness is validated against NetworkX and a Bellman optimality certificate in $O(|V|+|E|)$.

The results do not support a universal winner. The 4-ary heap obtains the lowest median in 7 of 10 synthetic scales, with a median ratio of 0.96 relative to the binary heap in the sparse profile and 0.95 in the dense profile, but the interquartile ranges overlap in all 10 scales; the binary heap wins in 3 scales and in the OSM query (15.65 ms versus 19.87 ms for the 4-ary heap and 26.09 ms for the Fibonacci heap). The Fibonacci heap never obtains the lowest median and is 1.05 to 1.59 times slower than the binary heap. The $A^*$ extension reduces expansions from 887 to 823 ($-7.2\%$) without changing the optimal cost. A uniform weather factor of 1.0691 increases the estimated time by $6.91\%$ but does not change the route because it scales every weight equally. The central conclusion is that queue choice depends on scale, density, and runtime environment, and that the asymptotic advantage of the Fibonacci heap does not translate into a practical advantage in this Python implementation.

\vspace{0.4cm}
\noindent\textbf{Keywords:} Dijkstra's algorithm, priority queues, binary heap, 4-ary heap, Fibonacci heap, shortest paths, OpenStreetMap, Cusco, $A^*$, experimental analysis.
